%/*% pdflatex -file-line-error -interaction nonstopmode %

\documentclass{article}
\usepackage{graphicx}
\usepackage{amsmath}
\usepackage{amsfonts}
\usepackage{bookmark}
\usepackage{booktabs}
\usepackage{exercises}
\graphicspath{ {./images/test} }
\usepackage{indentfirst}
\usepackage{amsthm}
\usepackage{yhmath}
\usepackage{placeins}
\newcommand*{\oldneg}{\mathord{\sim}}
\newcommand\showdiv[1]{\overline{\smash{)}#1}}
\newcommand\Mydiv[1]{%
\smash{\raise.3ex\hbox{$\big)$}}\mkern-8mu
        \overline{\enspace\strut#1}}
\setlength{\tabcolsep}{5pt}
\theoremstyle{definition}
\newtheorem{definition}{Definition}

\setlength{\parindent}{.5in}
\setlength{\parskip}{12pt}
\renewcommand{\arraystretch}{2}

\title{3.6 Answers}
\begin{document}
\maketitle

\section*{1}
When you solve an equation algebraically, an apparent solution that must be rejected because it does not satisfy the original equation is called 
an extraneous solution.

\section*{2}
He could have first isolated the radicals on separate sides.

\section*{10}
\[ -4\sqrt{x}-6=-20 \] \[ -4\sqrt{x}=-14 \] \[ \sqrt{x}=\frac{7}{2} \] \[ x=\frac{49}{4} \]

\section*{11}
\[ \sqrt{-2x+3}-2=10 \] \[ \sqrt{-2x+3}=12 \] \[ -2x+3=144 \] \[ x=-\frac{141}{2} \]

\section*{12}
C

\section*{20}
\[ \sqrt[3]{4x+2}-6=-10 \] \[ \sqrt[3]{4x+2}=-4 \] \[ 4x+2=-64 \] \[ x=-\frac{33}{2} \]

\section*{21}
\[ -4\sqrt[3]{x+10}+3=15 \] \[ \sqrt[3]{x+10}=-3 \] \[ x+10=-27 \] \[ x=-37 \]

\section*{22}
We know that $\sqrt[3]{-3a+b}=c$. So, we let $c=3$, meaning $-3a+b=27$. We let $b=9$, meaning $-3a+9=27$ and $a=-6$. Thus, one such equation is 
\[ \sqrt[3]{-6x+9}=3 \]

\section*{30}
\[ (5x-19)^\frac{5}{6}=32 \] \[ 5x-19=32^\frac{6}{5} \] \[ 5x=32^\frac{6}{5}+19 \] \[ 5x=(2^5)^\frac{6}{5}+19 \] \[ 5x=2^6+19 \] \[ x=\frac{83}{5} \]

\section*{31}
\[ (3x+43)^\frac{2}{3}+22=38 \] \[ (3x+43)^\frac{2}{3}=16 \] \[ 3x+43=16^\frac{3}{2} \] \[ 3x+43=(4^2)^\frac{3}{2} \] \[ 3x+43=4^3 \]
\[ 3x=21 \] \[ x=7 \]

\section*{33}
The error here is that on step two, the left side of the equation was squared, but not the right side of the equation. Here is a correct solution:
\[ (x+7)^\frac{1}{2}=5 \] \[ x+7=25 \] \[ x=18 \]

\section*{40}
\[ x-2=\sqrt{\frac{3}{2}x-2} \] \[ (x-2)^2=\frac{3}{2}x-2 \] \[ x^2-4x+4=\frac{3}{2}x-2 \] \[ x^2-\frac{11}{2}x=-6 \]
\[ x^2-\frac{11}{2}x+(\frac{11}{4})^2=-6+(\frac{11}{4})^2 \textrm{ (Complete the square)} \] 
\[ (x-\frac{11}{4})^2=1.5625 \] \[ x-\frac{11}{4}=\pm\sqrt{1.5625} \]
\[ x=\frac{11}{4}\pm 1.25 \] \[ x=\frac{3}{2}\vee x=4 \]
Now we check for extraneous solutions:
\[ \frac{3}{2}-2=\sqrt{\frac{3}{2}\frac{3}{2}-2} \] \[ -\frac{1}{2}=\sqrt{\frac{1}{4}} \] \[ -\frac{1}{2}=\frac{1}{2} \]
and \[ 4-2=\sqrt{\frac{3}{2}(4)-2} \] \[ 2=2 \] meaning the only solution is $x=2$.

\section*{41}
\[ \sqrt[4]{3-8x^2}=2x \] \[ 3-8x^2=16x^4 \] \[ 16x^4+8x^2-3=0 \] \[ (4x^2-1)(4x^2+3)=0 \] 
\[ x=\pm\frac{1}{2}\vee x=\pm\sqrt{-\frac{3}{4}}=\pm\frac{\sqrt{3}}{2}i \] 
We check for extraneous solutions:
\[ \sqrt[4]{3-8(\frac{1}{2})^2}=2(\frac{1}{2}) \] \[ 1=1 \] and \[ \sqrt[4]{3-8(-\frac{1}{2})^2}=2(-\frac{1}{2}) \]
\[ 1=-1 \]
and \[ \sqrt[4]{3-8(\frac{\sqrt{3}}{2}i)^2}=2(\frac{\sqrt{3}}{2}i) \] \[ \sqrt[4]{3-8(-\frac{3}{4})}=\sqrt{3}i \]
\[ \sqrt[4]{9}=\sqrt{3}i \] and \[ \sqrt[4]{3-8(-\frac{\sqrt{3}}{2}i)^2}=2(-\frac{\sqrt{3}}{2}i) \]
\[ \sqrt[4]{3-8(\frac{3}{4})}=-\sqrt{3}i \] \[ \sqrt[4]{-3}=-\sqrt{3}i \] \[ \sqrt[4]{3}i=-\sqrt{3}i \]
meaning $x=\frac{1}{2}$ is the only solution.

\section*{42}
\[ \sqrt[3]{8x^3-1}=2x-1 \] \[ 8x^3-1=(2x-1)^3 \] \[ 8x^3-1=8x^3-12x^2+6x-1 \] \[ 12x^2-6x=0 \] \[ x(12x-6)=0 \] \[ x=0\vee x=\frac{1}{2} \]
We check for extraneous solutions:
\[ \sqrt[3]{8(0)^3-1}=2(0)-1 \] \[ \sqrt[3]{-1}=-1 \] \[ -1=-1 \]
and \[ \sqrt[3]{8(\frac{1}{2})^2-1}=2(\frac{1}{2})-1 \] \[ 1=0 \] meaning $x=0$ is the only solution.

\section*{43}
\[ \sqrt{32x-64}=2x \] \[ 32x-64=4x^2 \] \[ 4x^2-32x+64=0 \] \[ 4(x^2-8+16)=0 \] \[ (x-4)^2=0 \] \[ x=4 \]
A

\section*{51}
\[ \sqrt{2x+5}=\sqrt{x+2}+1 \] \[ 2x+5=(\sqrt{x+2}+1)^2 \] \[ 2x+5=x+3+2\sqrt{x+2} \] \[ x+2=2\sqrt{x+2} \] \[ (x+2)^2=4x+8 \]
\[ x^2+4x+4=4x+8 \] \[ x^2-4=0 \] \[ x=\pm 2 \]
We check for extraneous solutions:
\[ \sqrt{2(2)+5}=\sqrt{2+2}+1 \] \[ 3=3 \] and \[ \sqrt{2(-2)+5}=\sqrt{-2+2}+1 \] \[ 1=1 \]
meaning both of $\pm 2$ are solutions. 

\section*{52}
\[ \sqrt{5x+6}+3=\sqrt{3x+3}+4 \] \[ \sqrt{5x+6}=\sqrt{3x+3}+1 \] \[ 5x+6=(\sqrt{3x+3}+1)^2 \] \[ 5x+6=3x+4+2\sqrt{3x+3} \]
\[ 2x+2=2\sqrt{3x+3} \] \[ (2x+2)^2=4(3x+3) \] \[ 4x^2+8x+4=12x+12 \] \[ 4x^2-4x-8=0 \] \[ 4(x^2-x-2)=0 \] \[ (x-2)(x+1)=0 \] \[ x=2\vee x=-1 \]
We check for extraneous solutions:
\[ \sqrt{5(2)+6}+3=\sqrt{3(2)+3}+4 \] \[ 7=7 \] and \[ \sqrt{5(-1)+6}+3=\sqrt{3(-1)+3}+4 \] \[ 4=4 \] meaning both $2$ and $-1$ are solutions. 

\section*{54}
First we solve for $x$ in the first equation:
\[ 5\sqrt{x}-2\sqrt{y}=4\sqrt{2} \] \[ 5\sqrt{x}=2\sqrt{y}+4\sqrt{2} \] \[ x=\frac{(2\sqrt{y}+4\sqrt{2})^2}{25} \]
Next, we plug this equation for $x$ into the second equation:
\[ 2\sqrt{\frac{(2\sqrt{y}+4\sqrt{2})^2}{25}}+3\sqrt{y}=13\sqrt{2} \]
\[ 2(\frac{2\sqrt{y}+4\sqrt{2}}{5})+3\sqrt{y}=13\sqrt{2} \]
\[ 2(2\sqrt{y}+4\sqrt{2})=-15\sqrt{y}+65\sqrt{2} \]
\[ 4\sqrt{y}+8\sqrt{2}=-15\sqrt{y}+65\sqrt{2} \]
\[ 19\sqrt{y}=57\sqrt{2} \]
\[ 361y=6498 \]
\[ y=18 \]
We then plug this value for $y$ into the first equation:
\[ 5\sqrt{x}-2\sqrt{18}=4\sqrt{2} \] \[ 5\sqrt{x}-6\sqrt{2}=4\sqrt{2} \] \[ 5\sqrt{x}=10\sqrt{2} \] \[ \sqrt{x}=2\sqrt{2} \] \[ x=8 \]
Thus the solution to this system of equations is $x=8,y=18$.

\section*{61}
(a): \[ 0=1.69\sqrt{s+4.25}-3.55 \] \[ \frac{3.55}{1.69}=\sqrt{s+4.25} \] \[ 4.41=s+4.25 \] \[ s=0.16 \] miles per hour.\newline
(b): \[ 12=1.69\sqrt{s+4.25}-3.55 \] \[ \frac{15.55}{1.69}=\sqrt{s+4.25} \] \[ 84.66=s+4.25 \] \[ s=80.41 \] miles per hour.\newline
(c): We know that the Beaufort numbers range from 0 to 12, meaning the wind speeds range from the speeds for the Beauforts numbers from 0 to 12. 
Because we know the wind speeds when $B=0$ and $B=12$, we know that the wind speed $s$ satisfies the inequality \[ 0.16\le s\le 80.41 \]

\section*{62}
\[ 5=\sqrt{h^2+\frac{1}{4}(4-2)^2} \] \[ 5=\sqrt{h^2+1} \] \[ 25=h^2+1 \] \[ h=\sqrt{24} \] \[ h=2\sqrt{6} \]

\end{document}
